% ===================================================================
% preamble.tex
% Configuraciones globales, paquetes y comandos personalizados
% ===================================================================

% ===================================================================
% PAQUETES Y CONFIGURACIONES DE IDIOMA
% ===================================================================
\usepackage[utf8]{inputenc}
\usepackage[spanish]{babel}
\usepackage{csquotes}
\usepackage{graphicx}
\usepackage{booktabs}
\usepackage{array}
\usepackage{multirow}
\usepackage{colortbl}
\usepackage{tikz}
\usepackage{amsmath}
\usepackage{amsfonts}
\usepackage{amssymb}
\usepackage{hyperref}
\usepackage{xcolor}
\usepackage{listings}
\usepackage{caption}
\usepackage{subcaption}
\usepackage{siunitx}
\usepackage{pgfplots}
\pgfplotsset{compat=1.17}
\usetikzlibrary{shapes,arrows,positioning,fit,backgrounds}

% ===================================================================
% CONFIGURACIONES DE BEAMER
% ===================================================================
\usetheme{Madrid}
\usecolortheme{default}

% Configuración de colores personalizados
\definecolor{primary}{RGB}{0, 51, 102}      % Azul institucional
\definecolor{secondary}{RGB}{184, 1, 68}    % Rojo corporativo
\definecolor{accent}{RGB}{46, 134, 171}     % Azul acento
\definecolor{success}{RGB}{46, 204, 113}    % Verde éxito
\definecolor{warning}{RGB}{241, 196, 15}    % Amarillo advertencia
\definecolor{danger}{RGB}{231, 76, 60}      % Rojo peligro

% Configuración del tema de colores
\setbeamercolor{title}{fg=white, bg=primary}
\setbeamercolor{frametitle}{fg=white, bg=primary}
\setbeamercolor{structure}{fg=primary}

% Personalización de bloques
\setbeamercolor{block title}{bg=primary, fg=white}
\setbeamercolor{block body}{bg=primary!10}

\setbeamercolor{block title example}{bg=success, fg=white}
\setbeamercolor{block body example}{bg=success!10}

\setbeamercolor{block title alerted}{bg=danger, fg=white}
\setbeamercolor{block body alerted}{bg=danger!10}

% Configuración de la tabla de contenidos
\setbeamertemplate{section in toc}{%
    \hspace{0.5em}%
    \inserttocsectionnumber.~\inserttocsection%
}

% Configuración de los footers
\setbeamertemplate{footline}{%
    \leavevmode%
    \hbox{%
        \begin{beamercolorbox}[wd=.333333\paperwidth,ht=2.25ex,dp=1ex,center]{author in head/foot}%
            \usebeamerfont{author in head/foot}\insertshortauthor
        \end{beamercolorbox}%
        \begin{beamercolorbox}[wd=.333333\paperwidth,ht=2.25ex,dp=1ex,center]{title in head/foot}%
            \usebeamerfont{title in head/foot}\insertshorttitle
        \end{beamercolorbox}%
        \begin{beamercolorbox}[wd=.333333\paperwidth,ht=2.25ex,dp=1ex,right]{date in head/foot}%
            \usebeamerfont{date in head/foot}\insertshortdate{}\hspace*{2em}
            \insertframenumber{} / \inserttotalframenumber\hspace*{2ex}
        \end{beamercolorbox}%
    }%
    \vskip0pt%
}

% ===================================================================
% COMANDOS PERSONALIZADOS
% ===================================================================

% Variables del proyecto
\newcommand{\churnrate}{26.5\%}
\newcommand{\datasetname}{Telco Customer Churn}
\newcommand{\projecttitle}{Análisis Predictivo de Churn en Telecomunicaciones}
\newcommand{\projectauthor}{Equipo de Business Intelligence}
\newcommand{\projectdate}{\today}

% Información de la presentación
\title[\projecttitle]{\projecttitle}
\subtitle{Presentación de Resultados y Propuesta de Acción}
\author[\projectauthor]{\projectauthor}
\institute[Telecom Corp]{
    \textbf{Telecom Corporation S.A.} \\
    \textit{Departamento de Business Intelligence}
}
\date{\projectdate}

% ===================================================================
% ENTORNOS PERSONALIZADOS
% ===================================================================

% Bloque de pregunta clave
\newenvironment{questionblock}{%
    \setbeamercolor{block title}{bg=secondary, fg=white}
    \setbeamercolor{block body}{bg=secondary!10}
    \begin{block}{\textbf{Pregunta Clave}}
}{\end{block}}

% Bloque de acción recomendada
\newenvironment{actionblock}{%
    \setbeamercolor{block title}{bg=primary, fg=white}
    \setbeamercolor{block body}{bg=primary!10}
    \begin{block}{\textbf{Acción Recomendada}}
}{\end{block}}

% Bloque de insight
\newenvironment{insightblock}{%
    \setbeamercolor{block title}{bg=accent, fg=white}
    \setbeamercolor{block body}{bg=accent!10}
    \begin{block}{\textbf{Insight Clave}}
}{\end{block}}

% Bloque de métrica
\newenvironment{metricblock}{%
    \setbeamercolor{block title}{bg=success, fg=white}
    \setbeamercolor{block body}{bg=success!10}
    \begin{block}{\textbf{Métrica}}
}{\end{block}}

% ===================================================================
% CONFIGURACIONES DE GRÁFICOS Y FIGURAS
% ===================================================================

% Ruta de imágenes
\graphicspath{{assets/images/}}

% Configuración de tablas
\newcolumntype{L}[1]{>{\raggedright\let\newline\\\arraybackslash\hspace{0pt}}m{#1}}
\newcolumntype{C}[1]{>{\centering\let\newline\\\arraybackslash\hspace{0pt}}m{#1}}
\newcolumntype{R}[1]{>{\raggedleft\let\newline\\\arraybackslash\hspace{0pt}}m{#1}}

% ===================================================================
% CONFIGURACIONES DE LISTADOS DE CÓDIGO
% ===================================================================

\lstset{
    basicstyle=\ttfamily\small,
    keywordstyle=\color{primary}\bfseries,
    commentstyle=\color{accent}\itshape,
    stringstyle=\color{secondary},
    showstringspaces=false,
    breaklines=true,
    frame=single,
    backgroundcolor=\color{gray!5},
    tabsize=2
}

% ===================================================================
% COMANDOS ÚTILES PARA PRESENTACIONES
% ===================================================================

% Resaltar texto
\newcommand{\highlight}[1]{{\color{secondary}\textbf{#1}}}
\newcommand{\highlightprimary}[1]{{\color{primary}\textbf{#1}}}
\newcommand{\highlightaccent}[1]{{\color{accent}\textbf{#1}}}

% Espaciado vertical
\newcommand{\vsp}{\vspace{0.5cm}}
\newcommand{\vspb}{\vspace{1cm}}

% Indicador de progreso
\newcommand{\progress}{%
    \begin{tikzpicture}[remember picture, overlay]
        \draw[primary!20] (current page.south west) rectangle (current page.south east);
        \fill[primary] (current page.south west) rectangle ++(\insertframenumber/\inserttotalframenumber\paperwidth, 3pt);
    \end{tikzpicture}%
}